% FaultRay: In-Memory Infrastructure Resilience Simulation
% arXiv submission — cs.SE / cs.DC
% Compile: pdflatex faultray-paper.tex && bibtex faultray-paper && pdflatex faultray-paper.tex && pdflatex faultray-paper.tex

\documentclass[10pt,twocolumn]{article}

%% ---- packages ----
\usepackage[utf8]{inputenc}
\usepackage[T1]{fontenc}
\usepackage{amsmath,amssymb,amsthm}
\usepackage{graphicx}
\usepackage{booktabs}
\usepackage{hyperref}
\usepackage[ruled,vlined,linesnumbered]{algorithm2e}
\usepackage{url}
\usepackage{xcolor}
\usepackage{balance}
\usepackage[margin=0.75in]{geometry}
\usepackage{enumitem}
\usepackage{multirow}
\usepackage{tabularx}

%% ---- theorem environments ----
\newtheorem{theorem}{Theorem}
\newtheorem{lemma}[theorem]{Lemma}
\newtheorem{corollary}[theorem]{Corollary}
\newtheorem{definition}{Definition}
\newtheorem{property}{Property}

%% ---- metadata ----
\title{FaultRay: In-Memory Infrastructure Resilience Simulation\\with Formal Cascade Analysis\\and Multi-Layer Availability Limits}

\author{
  Yutaro Maeda\\
  Independent Researcher\\
  Chigasaki, Japan\\
  \texttt{contact@faultray.com}
}

\date{}

\begin{document}
\maketitle

%% ====================================================================
%%  ABSTRACT
%% ====================================================================
\begin{abstract}
Production fault injection tools (Gremlin, Steadybit, AWS~FIS) are
powerful but touch running systems, and classical reliability
analysis (Fault Tree Analysis, Reliability Block Diagrams) operates
on static trees under component independence assumptions that do
not hold for cloud infrastructure where a single underlying
incident can simultaneously degrade multiple categorical layers.
We present \textsc{FaultRay}, a research prototype that contributes
two core technical elements:
(1)~a cascade propagation engine formalized as a Labeled Transition
System (LTS) over a dependency graph, with termination,
$O(|V|+|E|)$ time complexity, monotonicity, and bounded blast radius
(proof sketches provided; not mechanically verified);
and (2)~an $N$-layer availability limit model that decomposes a
system's availability ceiling via min-composition across software,
hardware, runtime noise floor, operational, and external SLA layers,
explicitly departing from the independence assumption of classical
reliability block diagrams (RBD).
We verify the engine's internal consistency via a post-hoc backtest
on 36~real-world cloud incidents spanning 2017--2023 (AWS, GCP,
Azure, Meta, Cloudflare), where topologies were constructed from
incident reports and cascade path reproduction is expected by
construction (29 of 36 achieve $F_1 = 1.000$; 7 exhibit modeling
imperfections). We do \emph{not} claim prospective prediction
capability; perturbation analysis confirms the engine is
non-degenerate but does not establish predictive validity.
We acknowledge concurrent complementary work by Krasnovsky
(ICSE-NIER~2026, Distinguished Paper Award; arXiv:2506.11176), which uses Monte-Carlo graph
simulation without formal proofs or multi-layer decomposition.
\textsc{FaultRay} is released under Apache~License~2.0 on GitHub
and PyPI. A provisional patent application has been filed with the
USPTO (Application No.\ 64/010,200, 2026-03-19).

\textbf{This is a research prototype and NOT a certified compliance
tool. Outputs must not be used as audit evidence without independent
legal and technical review.}
\end{abstract}

%% ====================================================================
%%  1. INTRODUCTION
%% ====================================================================
\section{Introduction}\label{sec:intro}

The complexity of modern distributed infrastructure continues to grow.
A typical production system may involve load balancers, application
servers, databases, caches, message queues, DNS resolvers, and storage
tiers, all interconnected by dependency edges that amplify localized
failures into system-wide outages.  The 2021 Facebook (Meta) BGP
misconfiguration~\cite{meta2021}, the 2017 Amazon S3 key-space
corruption~\cite{aws-s3-2017}, and the 2022 Cloudflare BGP
leak~\cite{cloudflare2022} each demonstrate how a single fault can
cascade across dozens of dependent services within minutes.

Current chaos engineering approaches suffer from several limitations:
\begin{itemize}[nosep]
\item \textbf{Production risk.} Netflix Chaos Monkey~\cite{basiri2016},
  Gremlin~\cite{gremlin2020}, Steadybit~\cite{steadybit2023}, and
  AWS FIS~\cite{aws-fis} inject real faults into running systems. For
  regulated financial institutions and other risk-averse operators,
  touching production to discover weaknesses is not always acceptable.
\item \textbf{Static analysis without cascade semantics.}
  HPE's SPOF detection~\cite{hpe-spof} identifies single points of
  failure but does not model cascade propagation through typed
  dependencies.
\item \textbf{Ad hoc availability numbers.} Classical MTBF/MTTR
  models~\cite{trivedi2001} and Reliability Block Diagrams produce
  a single availability number under an independence assumption that
  does not match cloud infrastructure, where a single vendor incident
  can simultaneously degrade multiple layers (compute, network,
  storage) in correlated fashion.
\end{itemize}

The idea of simulating failures without touching production is not
new: digital twin approaches (ChaosTwin~\cite{chaostwin2021},
Farokhi et~al.~\cite{farokhi2024}) replicate full system behavior
in isolation, and concurrent work by
Krasnovsky~\cite{krasnovsky2026} (ICSE-NIER~2026) uses Monte-Carlo
sampling over trace-derived graphs. We view these as complementary
operating points (\S\ref{sec:rw-krasnovsky}, \S\ref{sec:rw-digital-twin}):
\textsc{FaultRay} trades fidelity for analytical tractability via
a static dependency graph with formal LTS semantics.

\noindent
This paper contributes two core technical elements:

\begin{enumerate}[nosep]
\item A \textbf{cascade propagation engine} formalized as a Labeled
  Transition System over a dependency graph, with termination
  on finite graphs, $O(|V| + |E|)$ time complexity,
  monotonicity under additional fail-stop events, causality
  preservation, and a bounded blast radius
  (\S\ref{sec:cascade}).
\item An \textbf{$N$-layer availability limit model} that decomposes
  a system's availability ceiling across software, hardware,
  runtime noise floor, operational, and external SLA layers using
  \emph{min-composition} rather than the classical series-product,
  and provides a weakest-link justification for that choice
  (\S\ref{sec:availability}).
\end{enumerate}

\textsc{FaultRay} is released as open source under Apache~License~2.0.
A provisional patent application covering the cascade engine and
$N$-layer decomposition has been filed with the USPTO (Application
No.~64/010,200, filed 2026-03-19). This paper is a research
prototype description; it is not a certified compliance tool and
its outputs must not be used as audit evidence without independent
legal and technical review.

%% ====================================================================
%%  2. RELATED WORK
%% ====================================================================
\section{Related Work}\label{sec:related}

\subsection{Chaos Engineering (Production Fault Injection)}
Netflix introduced Chaos Monkey in 2011, later formalized as
\emph{Chaos Engineering} by Basiri et al.~\cite{basiri2016}.
Commercial platforms (Gremlin, Steadybit, AWS~FIS) and open-source
CNCF projects (LitmusChaos, Chaos Mesh) provide managed fault
injection in production or staging environments; a comprehensive
survey is provided by Joshua et~al.~\cite{chaos-mlr-2025}.
All of these tools require running
infrastructure and inject \emph{real} faults, introducing production
risk. \textsc{FaultRay} is strictly analytical and never touches
running infrastructure, which is a distinct operating point.

\subsection{In-Memory Graph Simulation (Concurrent Work)}
\label{sec:rw-krasnovsky}
Concurrent and complementary work by
Krasnovsky~\cite{krasnovsky2026} (ICSE-NIER~2026, Distinguished Paper Award; arXiv:2506.11176)
proposes \emph{model discovery}: extracting a service-dependency
graph from trace data (e.g., Jaeger) and running Monte-Carlo
fail-stop simulation over it. Evaluated on DeathStarBench, the
approach achieves mean absolute error $\leq 0.0004$ against
observed resilience.
\textsc{FaultRay} and Krasnovsky~2026 share the high-level positioning
of lightweight, in-memory simulation as an alternative to production
fault injection. The approaches differ in three technical axes:
(a) Krasnovsky uses Monte-Carlo sampling, whereas \textsc{FaultRay}
uses a formal LTS with termination and $O(|V|+|E|)$ complexity
arguments (proof sketches, not mechanically verified) and bounded
blast radius;
(b) Krasnovsky treats the graph as a single-layer connectivity
model, whereas \textsc{FaultRay} decomposes availability across $N$
infrastructure layers with min-composition;
(c) Krasnovsky emphasizes trace-based auto-discovery of the graph,
whereas \textsc{FaultRay} accepts manually-specified or
infrastructure-as-code-derived graphs.
We view the two works as complementary and not substitutes.

\subsection{Digital Twins and In-Memory Resilience Testing}
\label{sec:rw-digital-twin}
The concept of testing system resilience without touching
production infrastructure has been explored in the digital twin
community. ChaosTwin~\cite{chaostwin2021} (CNSM~2021) combines
digital twin replicas with chaos engineering principles to test
resilience in an isolated environment, demonstrating the approach
on network function virtualization. Farokhi
et~al.~\cite{farokhi2024} (IEEE~TII~2024) extend this to
industrial digital twins, proposing chaos injection into the
digital replica rather than the physical system.
\textsc{FaultRay} shares the high-level positioning of
``simulate failures without production impact'' but differs in
method: digital twin approaches replicate the full system
behavior (including runtime dynamics), whereas \textsc{FaultRay}
operates on a static dependency graph with formal LTS semantics.
Digital twins are higher fidelity but require a running replica;
\textsc{FaultRay} sacrifices fidelity for analytical tractability
and formal guarantees.

A recent multi-vocal literature review on chaos
engineering~\cite{chaos-mlr-2025} (ACM~TOSEM~2025) synthesizes
both academic and industry perspectives. That survey positions
in-memory analytical simulation as a distinct research direction
from production fault injection, which is the operating point
\textsc{FaultRay} occupies.

\subsection{Formal Methods for Distributed Systems}
\label{sec:rw-formal}
Labeled Transition Systems have been used in formal verification
of distributed systems for decades, originating with
Keller~\cite{keller1976}. Recent work on Validating Labelled
State Transition and Message Production Systems
(VLSTS)~\cite{vlsts2022} proposes LTS-based modeling and
verification of distributed systems with faults, focusing on
message-passing correctness properties. \textsc{FaultRay}'s use
of LTS is conceptually simpler: transitions model health state
degradation rather than message-passing, and the proofs
(termination, monotonicity, blast radius) are correspondingly
straightforward. The novelty is not in the LTS formalism itself
but in its application to infrastructure dependency graphs with
typed dependencies and multi-layer availability composition.
More recently, Nasar~\cite{nasar2026} applied TLA+ to formally verify
microgrid reliability models, demonstrating that mechanical
verification of infrastructure availability properties is feasible
with current tools. \textsc{FaultRay}'s proof sketches have not been
mechanically verified (\S\ref{sec:discussion}); we consider this an
important gap.

\subsection{Static Resilience Analysis}
HPE's single-point-of-failure detection system
(US Patent 9,280,409~B2)~\cite{hpe-spof} analyzes dependency graphs
to identify components whose failure would cause total outage.
However, it performs \emph{static} analysis---it does not simulate
cascade propagation through dependency types (required, optional,
asynchronous) or account for circuit breakers, retry strategies, or
latency amplification.

\subsection{Classical Reliability Engineering (FTA and RBD)}
\label{sec:rw-fta-rbd}
Fault Tree Analysis (FTA), standardized in NASA NUREG-0492~\cite{nasa-fta},
IEC~61025~\cite{iec61025}, and SAE~ARP4761~\cite{arp4761}, constructs
a static AND/OR gate tree and enumerates minimal cut sets to compute
failure probabilities. Reliability Block Diagrams (RBD), treated
exhaustively by Trivedi~\cite{trivedi2001}, compose component
availabilities using series-product
($A_{\text{series}} = \prod_i A_i$) and parallel-redundancy
($A_{\text{parallel}} = 1 - \prod_i (1 - A_i)$) operators under an
independence assumption.
\textsc{FaultRay} departs from FTA and RBD in two ways:
(i)~\textsc{FaultRay}'s cascade engine is a \emph{dynamic} LTS over
a directly-specified dependency graph rather than a hand-constructed
AND/OR tree;
(ii)~\textsc{FaultRay}'s $N$-layer availability model uses
\emph{min-composition} rather than series-product composition, which
is appropriate when infrastructure layers are not independent
(e.g., a cloud provider outage can simultaneously degrade compute,
network, and storage layers in correlated fashion).

\subsection{SLA Decomposition}
\label{sec:rw-sla}
Casformer~\cite{casformer2025} (arXiv:2601.11859) uses a cascaded
Transformer architecture to learn end-to-end SLA decomposition for
6G network slicing. Broader surveys on cloud SLA
decomposition~\cite{sla-survey-2024} discuss iterative optimization
approaches. These works target telecommunications and network
slicing, and use learned or optimization-based aggregation,
whereas \textsc{FaultRay} targets cloud software infrastructure
availability with formal min-composition.

\subsection{Cascade Failures in Interdependent Networks}
Buldyrev et al.~\cite{buldyrev2010} demonstrated that
interdependent networks are dramatically more fragile than
isolated networks, establishing the theoretical foundation for
cascade failure modeling across coupled systems. The I$^3$
model~\cite{iii-model-2025} (arXiv:2503.02890) extends this to
urban infrastructure using graph autoencoders. Broader work on
cascading failures is surveyed in~\cite{pnas-cascade-2019}. The
fundamental taxonomy of faults, errors, and failures was
established by Avizienis et al.~\cite{avizienis2004}; \textsc{FaultRay}'s
health status model (\textsc{Healthy}/\textsc{Degraded}/\textsc{Overloaded}/\textsc{Down})
is a simplification of this taxonomy. These foundational works
target physical or abstract systems; \textsc{FaultRay} applies
similar concepts specifically to cloud software dependency graphs.

\subsection{Availability Modeling}
Classical availability theory uses the formula
$A = \text{MTBF} / (\text{MTBF} + \text{MTTR})$~\cite{trivedi2001},
extended to parallel redundancy via
$A_{\text{tier}} = 1 - (1 - A_{\text{single}})^{n}$.
In practice, availability is limited by factors beyond hardware
reliability---deployment frequency, human error, operational
response time, and external SLA dependencies.
\textsc{FaultRay}'s $N$-layer decomposition is an attempt to
expose these factors as separately-analyzable layers whose
effective ceiling is determined by the weakest link (min
composition).

\subsection{Positioning Summary}
\label{sec:rw-gap}
\textsc{FaultRay}'s contribution is not a single novel idea but a
specific combination: formal LTS cascade semantics with proven
termination and complexity bounds (differentiating from
Krasnovsky~2026 and production chaos tools), applied to software
dependency graphs (differentiating from FTA/RBD hardware
reliability), with min-composition multi-layer availability
(differentiating from RBD series-product and Casformer learned
aggregation). Table~\ref{tab:gap} summarizes the positioning.

\begin{table}[t]
\centering
\caption{Positioning of \textsc{FaultRay} relative to representative
prior work. ``In-mem'' = no production fault injection.
``Formal'' = termination and complexity arguments (proof sketches).
``N-layer'' = multi-layer availability decomposition.
``OSS'' = open source release.}
\label{tab:gap}
\small
\begin{tabular}{lcccc}
\toprule
\textbf{Tool / Work} & \textbf{In-mem} & \textbf{Formal} & \textbf{N-layer} & \textbf{OSS} \\
\midrule
Chaos Monkey / Gremlin / Steadybit / AWS FIS & \texttimes & \texttimes & \texttimes & partial \\
ChaosTwin (CNSM~2021) & \checkmark\textsuperscript{*} & \texttimes & \texttimes & \texttimes \\
Farokhi~2024 (IEEE~TII) & \checkmark\textsuperscript{*} & \texttimes & \texttimes & \texttimes \\
HPE SPOF (US~9,280,409) & \checkmark & \texttimes & \texttimes & \texttimes \\
FTA (NASA/IEC/SAE) & \checkmark & partial & hier.\ AND/OR & \texttimes \\
RBD (Trivedi~2001) & \checkmark & \texttimes & series/parallel & \texttimes \\
VLSTS (arXiv:2202.12662) & \checkmark & \checkmark & \texttimes & \texttimes \\
Krasnovsky~2026 (arXiv:2506.11176) & \checkmark & \texttimes & \texttimes & \checkmark \\
Casformer (arXiv:2601.11859) & N/A & \texttimes & learned (6G) & \texttimes \\
I$^3$ model (arXiv:2503.02890) & \checkmark & \texttimes & \texttimes & partial \\
\textbf{\textsc{FaultRay}} & \checkmark & \checkmark & min-composition & \checkmark \\
\bottomrule
\end{tabular}
\\[2pt]
{\footnotesize \textsuperscript{*}Digital twin approaches replicate
  the full system rather than using a static dependency graph.}
\end{table}

%% ====================================================================
%%  3. SYSTEM ARCHITECTURE
%% ====================================================================
\section{System Architecture}\label{sec:architecture}

\subsection{Graph-Based Topology Model}
\textsc{FaultRay} represents infrastructure as a directed graph
$G = (C, E)$ where each vertex $c \in C$ is a typed component and
each edge $e \in E$ is a typed dependency.

\begin{definition}[Infrastructure Graph]\label{def:graph}
An infrastructure graph is a tuple $G = (C, E, \mathit{dep}, w, \tau, \mathit{cb})$
where:
\begin{itemize}[nosep]
\item $C$: finite set of components, each with type
  $\mathit{type}(c) \in \{\texttt{load\_balancer}, \texttt{web\_server},
  \texttt{app\_server}, \texttt{database}, \texttt{cache}, \texttt{queue},
  \texttt{storage}, \texttt{dns}, \texttt{external\_api},
  \texttt{custom}\}$;
\item $E \subseteq C \times C$: directed dependency edges;
\item $\mathit{dep}: E \to \{\texttt{requires}, \texttt{optional}, \texttt{async}\}$:
  dependency type;
\item $w: E \to [0, 1]$: dependency weight (criticality);
\item $\tau: E \to \mathbb{R}_{\ge 0}$: edge latency (ms);
\item $\mathit{cb}: E \to \{\mathit{true}, \mathit{false}\}$: circuit breaker status.
\end{itemize}
\end{definition}

Each component $c$ carries capacity metadata (max connections, max RPS,
timeout, retry multiplier, connection pool size), resource metrics
(CPU, memory, disk, network connections), and redundancy configuration
(replicas, autoscaling, failover). The open-source implementation
additionally supports application-specific component types (including
LLM-based agents) as a tool capability; such extensions are not
claimed as a research contribution in this paper.

\subsection{Automated Fault Scenario Generation}
\textsc{FaultRay} automatically generates fault scenarios from the
topology graph.  Given $|C|$ components and a fault type catalog of
$k$~types, the engine produces $|C| \times k$ single-fault scenarios
plus ${|C| \choose 2} \times k^2$ correlated-fault scenarios.
The core fault types covered by this paper are infrastructure-level
fail-stop events and capacity degradations. Each scenario is a tuple
$\langle \mathit{target}, \mathit{type}, \mathit{severity},
\mathit{duration} \rangle$.

\begin{figure}[t]
\centering
\small
\setlength{\unitlength}{1mm}
\begin{picture}(75,52)
\put(0,0){\framebox(75,52){}}
\put(22,46){\makebox(30,4){\textbf{FaultRay Engine}}}
\put(2,34){\framebox(20,9){\shortstack{Topology\\Parser}}}
\put(27,34){\framebox(20,9){\shortstack{Scenario\\Generator}}}
\put(52,34){\framebox(20,9){\shortstack{Avail.\\Model}}}
\put(22,38){\vector(1,0){5}}
\put(47,38){\vector(1,0){5}}
\put(37,34){\vector(0,-1){4}}
\put(27,21){\framebox(20,9){\shortstack{Cascade\\Engine (LTS)}}}
\put(37,21){\vector(0,-1){4}}
\put(62,34){\vector(0,-1){21}}
\put(27,8){\framebox(20,9){\shortstack{Backtest\\Harness}}}
\put(37,8){\vector(0,-1){4}}
\put(5,0){\framebox(65,4){\small Report / API / CLI}}
\end{picture}
\caption{System architecture of \textsc{FaultRay}. The cascade engine
(formalized as an LTS, \S\ref{sec:cascade}) and the $N$-layer
availability model (\S\ref{sec:availability}) are the central
components; the backtest harness and reporting layer are supporting
infrastructure.}
\label{fig:architecture}
\end{figure}

%% ====================================================================
%%  4. CASCADE ENGINE — FORMAL SPECIFICATION
%% ====================================================================
\section{Cascade Engine --- Formal Specification}\label{sec:cascade}

\subsection{Labeled Transition System}

\begin{definition}[Cascade Propagation Semantics]\label{def:lts}
The Cascade Propagation Semantics (CPS) is a Labeled Transition System
$\mathcal{M} = (S, s_0, \mathit{Act}, \longrightarrow, F)$ where:
\begin{itemize}[nosep]
\item $S$: set of states, each a 4-tuple $s = (H, L, T, V)$;
\item $s_0 = (H_0, L_0, 0, \emptyset)$: initial state;
\item $\mathit{Act}$: set of actions (fault injection, propagation, circuit breaker trip, timeout, termination);
\item ${\longrightarrow} \subseteq S \times \mathit{Act} \times S$: transition relation;
\item $F \subseteq S$: terminal states where no further transitions are enabled.
\end{itemize}
\end{definition}

\begin{definition}[State 4-Tuple]\label{def:state}
A CPS state $s = (H, L, T, V)$ consists of:
\begin{itemize}[nosep]
\item $H: C \to \{\textsc{Healthy}, \textsc{Degraded}, \textsc{Overloaded}, \textsc{Down}\}$
  --- health map;
\item $L: C \to \mathbb{R}_{\ge 0}$ --- accumulated latency map (ms);
\item $T \in \mathbb{R}_{\ge 0}$ --- elapsed simulation time (s);
\item $V \subseteq C$ --- visited set (monotonically growing).
\end{itemize}
Health status is totally ordered:
$\textsc{Healthy} < \textsc{Degraded} < \textsc{Overloaded} < \textsc{Down}$.
\end{definition}

\subsection{Transition Rules}

We define eight transition rules.  In each rule, unprimed variables
denote the pre-state and primed variables denote the post-state.
We write $H[c \mapsto h]$ for the map $H$ updated at component~$c$
to health~$h$.

\medskip\noindent
\textbf{Rule 1 (Fault Injection).}
Given fault $f = (\mathit{target}, \mathit{type}, \mathit{severity})$:
\begin{equation}
  \frac{c = f.\mathit{target} \quad h = \mathit{effect}(f.\mathit{type})}
  {(H, L, T, V) \xrightarrow{\textit{inject}(f)}
   (H[c \mapsto h], L, T, V \cup \{c\})}
\end{equation}
where $\mathit{effect}: \mathit{FaultType} \to \mathit{HealthStatus}$ maps
each fault type to the resulting health status (e.g.,
$\mathit{effect}(\texttt{component\_down}) = \textsc{Down}$,
$\mathit{effect}(\texttt{latency\_spike}) = \textsc{Degraded}$).

\medskip\noindent
\textbf{Rule 2 (Required Dependency --- Single Replica).}
If component $c'$ has a \texttt{requires} dependency on failed component~$c$,
$c$~has a single replica, and $c' \notin V$:
\begin{equation}
  \frac{\mathit{dep}(c', c) = \texttt{requires} \quad
        \mathit{replicas}(c) = 1 \quad
        H(c) = \textsc{Down}}
  {(H, L, T, V) \xrightarrow{\textit{prop}(c, c')}
   (H[c' \mapsto \textsc{Down}], L', T, V \cup \{c'\})}
\end{equation}

\medskip\noindent
\textbf{Rule 3 (Required Dependency --- Multiple Replicas).}
If the failed component $c$ has $n > 1$ replicas, its failure reduces
capacity but does not cause immediate failure of dependents:
\begin{equation}
  \frac{\mathit{dep}(c', c) = \texttt{requires} \quad
        \mathit{replicas}(c) > 1 \quad
        H(c) = \textsc{Down}}
  {(H, L, T, V) \xrightarrow{\textit{prop}(c, c')}
   (H[c' \mapsto \max(\textsc{Degraded}, H(c'))], L', T, V \cup \{c'\})}
\end{equation}
The $\max$ preserves monotonicity: if $c'$ is already in a state
worse than \textsc{Degraded} (due to a prior cascade path), that
state is retained.

\medskip\noindent
\textbf{Rule 4 (Optional Dependency).}
If $c'$ has an \texttt{optional} dependency on~$c$, the cascade is
attenuated regardless of $c$'s health:
\begin{equation}
  \frac{\mathit{dep}(c', c) = \texttt{optional} \quad
        H(c) \ge \textsc{Degraded}}
  {(H, L, T, V) \xrightarrow{\textit{prop}(c, c')}
   (H[c' \mapsto \max(\textsc{Degraded}, H(c'))], L, T, V \cup \{c'\})}
\end{equation}
The $\max$ ensures monotonicity: if $c'$ is \textsc{Healthy}, it
transitions to \textsc{Degraded}; if already worse, it is preserved.

\medskip\noindent
\textbf{Rule 5 (Async Dependency).}
If $c'$ has an \texttt{async} dependency on~$c$, the cascade is
attenuated (as with optional dependencies) but propagation is
\emph{delayed} by the edge latency $\tau(c', c)$:
\begin{equation}
  \frac{\mathit{dep}(c', c) = \texttt{async} \quad
        H(c) \ge \textsc{Degraded}}
  {(H, L, T, V) \xrightarrow{\textit{prop}(c, c')}
   (H[c' \mapsto \max(\textsc{Degraded}, H(c'))], L,
    T + \tau(c', c), V \cup \{c'\})}
\end{equation}
The delay models the asynchronous nature of the dependency:
message queues, event buses, and batch pipelines propagate
failures with latency rather than immediately. This distinguishes
\texttt{async} from \texttt{optional} semantically: both attenuate
health impact to at most \textsc{Degraded}, but \texttt{async}
additionally advances simulation time, affecting cascade chain
timestamps and severity scoring.

\medskip\noindent
\textbf{Rule 6 (Circuit Breaker Trip).}
If accumulated latency on the edge exceeds the timeout and the edge has
a circuit breaker enabled, the cascade is stopped at that edge:
\begin{equation}
  \frac{\mathit{cb}(c', c) = \mathit{true} \quad
        L(c) > \tau_{\max}(c')}
  {(H, L, T, V) \xrightarrow{\textit{cb\_trip}(c', c)}
   (H[c' \mapsto \max(\textsc{Degraded}, H(c'))], L, T, V \cup \{c'\})}
\end{equation}
where $\tau_{\max}(c')$ is the timeout threshold of component~$c'$.
The $\max$ operator preserves monotonicity
(Theorem~\ref{thm:monotonicity}): if $c'$ is already in a state
worse than \textsc{Degraded}, that state is retained.
The circuit breaker isolates $c'$ from further cascade propagation
through this edge.

\medskip\noindent
\textbf{Rule 7 (Timeout Propagation).}
If accumulated latency exceeds the timeout but no circuit breaker is
present, the component transitions to \textsc{Down} due to connection
pool exhaustion from retry storms:
\begin{equation}
  \frac{\mathit{cb}(c', c) = \mathit{false} \quad
        L(c') + \tau(c', c) > \tau_{\max}(c')}
  {(H, L, T, V) \xrightarrow{\textit{timeout}(c')}
   (H[c' \mapsto \textsc{Down}],\, L[c' \mapsto L(c') + \tau(c',c)],\, T,\, V \cup \{c'\})}
\end{equation}

\medskip\noindent
\textbf{Rule 8 (Termination).}
No transition is enabled when all reachable components have been visited:
\begin{equation}
  \frac{\forall c' \in \mathit{reach}(V).\; c' \in V}
  {(H, L, T, V) \in F}
\end{equation}

\subsection{Correctness Properties}

The following properties are stated with proof sketches.
These are \emph{informal} arguments intended to convey the
reasoning; they have not been mechanically verified in a
proof assistant (Coq, Isabelle, TLA+). We consider mechanical
verification an important future work item.

\begin{theorem}[Termination]\label{thm:termination}
For any infrastructure graph $G = (C, E, \ldots)$ and initial fault $f$,
the CPS $\mathcal{M}$ terminates in a finite number of transitions.
\end{theorem}
\begin{proof}[Proof sketch]
Define the progress measure $\phi(s) = |C| - |V|$ where $V$ is the
visited set in state~$s$. Every transition rule adds at least one
component to~$V$ (by inspection of Rules~1--7: each rule's consequent
includes $V \cup \{c'\}$), so $\phi$ strictly decreases on each
transition: $\phi(s') < \phi(s)$. Since $\phi$ is a natural number
bounded below by~0, the sequence of transitions is finite. The visited
set also prevents re-entry: a component $c' \in V$ cannot trigger a
new transition because all rules require processing a component not
yet in~$V$ (or, for cyclic graphs, the depth limit $D_{\max} = 20$
is enforced as an additional guard). The CPS therefore terminates in
at most $\min(|C|, D_{\max} \cdot \Delta)$ transitions, where
$\Delta$ is the maximum out-degree.
\end{proof}

\begin{theorem}[Time Complexity]\label{thm:complexity}
The CPS executes in $O(|C| + |E|)$ time.
\end{theorem}
\begin{proof}[Proof sketch]
Each component is visited at most once (the visited set~$V$ prevents
re-processing; proved in Theorem~\ref{thm:termination}). For each
visited component, its outgoing edges in the reverse dependency graph
are examined exactly once to identify dependents. The per-edge work
(health comparison, dependency type lookup) is $O(1)$. The total work
is therefore bounded by $\sum_{c \in C} (1 + \deg^{-}(c)) = |C| + |E|$.
\end{proof}

\begin{theorem}[Monotonicity of Failure]\label{thm:monotonicity}
Health can only worsen during a simulation run:
if $s \xrightarrow{*} s'$, then $\forall c \in C.\; H'(c) \ge H(c)$.
\end{theorem}
\begin{proof}[Proof sketch]
Define the health ordering \textsc{Healthy} $<$ \textsc{Degraded} $<$
\textsc{Overloaded} $<$ \textsc{Down}. We verify by case analysis on
each transition rule:
Rule~1 (Injection): sets $H(c_0) \leftarrow \mathit{effect}(f.\mathit{type})$.
Since $\mathit{effect}$ maps to a status $\ge \textsc{Degraded}$ and
the injection target is added to the visited set ($V \cup \{c_0\}$),
it will not be re-visited. For the initial injection,
$H(c_0) = \textsc{Healthy}$, so $H'(c_0) \ge H(c_0)$.
Rule~2 (Required, single replica): sets $H(c') \leftarrow \textsc{Down}$,
which is the maximum health status and therefore $\ge H(c')$ for any
prior health.
Rule~3 (Required, multiple replicas): sets
$H(c') \leftarrow \max(\textsc{Degraded}, H(c'))$, preserving
monotonicity by the $\max$ operator (identical argument to
Rules~4--5).
Rules~4--5 (Optional, Async): set $H(c') \leftarrow \max(\textsc{Degraded},
H(c'))$. Since $\max(\textsc{Degraded}, H(c')) \ge H(c')$ for any
health status, monotonicity is preserved: \textsc{Healthy} worsens to
\textsc{Degraded}, and states already worse than \textsc{Degraded} are
preserved.
Rule~6 (CB Trip): sets $H(c') \leftarrow \max(\textsc{Degraded}, H(c'))$,
which preserves monotonicity by the $\max$ operator (identical argument
to Rules~4--5).
Rule~7 (Timeout): sets $H(c') \leftarrow \textsc{Down}$, which is
$\ge H(c')$ for any prior health.
In all cases, $H'(c) \ge H(c)$.
\end{proof}

\begin{theorem}[Causality]\label{thm:causality}
For any component $c' \neq c_0$ (where $c_0$ is the initially
injected fault target), $c'$ transitions to a non-healthy state only
if at least one of its dependencies has already transitioned:
\begin{equation*}
  c' \neq c_0 \wedge H'(c') > \textsc{Healthy} \implies \exists c \in C.\;
  (c', c) \in E \wedge H(c) > \textsc{Healthy}
\end{equation*}
The initial fault target $c_0$ is an exception: it transitions via
Rule~1 (direct injection) without requiring a failed dependency.
\end{theorem}

\begin{theorem}[Circuit Breaker Correctness]\label{thm:cb}
If edge $(c', c)$ has $\mathit{cb}(c', c) = \mathit{true}$ and the
circuit breaker trips, then $c'$ transitions to
$\max(\textsc{Degraded}, H(c'))$---at least \textsc{Degraded},
and preserving any prior state that is already worse than
\textsc{Degraded}---and no further cascade propagation occurs
through this edge.
\end{theorem}

\begin{theorem}[Dependency Type Attenuation]\label{thm:attenuation}
Optional and asynchronous dependencies attenuate cascade effects:
if $\mathit{dep}(c', c) \in \{\texttt{optional}, \texttt{async}\}$
and $c'$ was previously \textsc{Healthy}, then
$H'(c') = \textsc{Degraded}$ regardless of $H(c)$. If $c'$ was
already in a state worse than \textsc{Degraded} (due to a required
dependency), that state is preserved by monotonicity
(Theorem~\ref{thm:monotonicity}). The cascade contribution of an
optional/async dependency is therefore bounded at \textsc{Degraded}
in \emph{severity}. Since \textsc{Degraded} does not trigger further
\texttt{requires}-dependency cascade to \textsc{Down}
(Rule~2 requires $H(c) = \textsc{Down}$), the
severity of cascade through optional/async paths is bounded.
However, the \emph{extent} (number of affected components) is bounded
only by the reverse-reachable set (Theorem~\ref{thm:blast}), as
\textsc{Degraded} propagation through chains of optional/async edges
can reach all transitively connected components.
\end{theorem}

\begin{theorem}[Blast Radius Bound]\label{thm:blast}
The blast radius of a single fault injection on component $c$ is
bounded by the set of components transitively reachable from $c$
via the reverse dependency graph:
\begin{equation*}
  \mathit{affected}(c) \subseteq \{c' \in C \mid c' \xrightarrow{*} c \text{ in } G^{-1}\}
\end{equation*}
Components not in the reverse-reachable set are guaranteed to be unaffected.
\end{theorem}

\begin{corollary}[Cyclic Graph Termination]\label{cor:cyclic}
For graphs with cycles, the depth limit $D_{\max} = 20$ together with
the visited set $V$ guarantees termination within at most
$\min(|C|, D_{\max} \cdot \Delta)$ transitions, where $\Delta$ is
the maximum out-degree.
\end{corollary}

\paragraph{Choice of $D_{\max} = 20$.}
The depth limit is a configurable engineering parameter, not a
theoretically derived value. We chose $D_{\max} = 20$ because
it exceeds the longest cascade chain observed in our 36-incident
test suite (maximum observed depth: 6) by a factor of~3, providing
headroom for larger topologies while bounding worst-case execution
in cyclic graphs. For very large service meshes ($>$500~components),
a practitioner may increase $D_{\max}$ or rely solely on the
visited set for termination. The formal termination guarantee
(Theorem~\ref{thm:termination}) holds for any finite $D_{\max}$.

%% ====================================================================
%%  5. N-LAYER AVAILABILITY LIMIT MODEL
%% ====================================================================
\section{$N$-Layer Availability Limit Model}\label{sec:availability}

Traditional availability is a single number:
$A = \text{MTBF} / (\text{MTBF} + \text{MTTR})$~\cite{trivedi2001}.
Classical Reliability Block Diagrams extend this with a
series-product composition rule
$A_{\text{series}} = \prod_i A_i$
under an assumption that components fail independently.
\textsc{FaultRay} takes a different position: in cloud infrastructure,
different categories of availability loss (deployment mistakes,
hardware faults, physical jitter, operational response time, external
SLA dependencies) are \emph{not} independent in the failure-causing
sense, because a single underlying incident---a provider region
outage, a kernel bug, a human deploy error---can simultaneously
degrade multiple categories.

We therefore decompose system availability into five
\emph{categorical layers}, each representing an upper bound imposed
by a distinct class of failure causes, and compose them with the
\emph{min} operator rather than the series-product:

\begin{definition}[System Availability Ceiling]
\begin{equation}
  A_{\text{system}} = \min(A_{\text{L1}},\, A_{\text{L2}},\, A_{\text{L3}},\, A_{\text{L4}},\, A_{\text{L5}})
\end{equation}
\end{definition}

\paragraph{Why min rather than product?}
There are three reasons to compose categorical availability ceilings
with $\min$ rather than $\prod$:

\emph{(1) Correlated failure modes.}
Classical RBD series-product
$\prod_i A_i$ assumes that component failures are
\emph{statistically independent}. For cloud infrastructure layers,
this assumption fails: a cloud provider region incident can
simultaneously reduce $A_{\text{L2}}$ (hardware availability as
observed through the VM layer), $A_{\text{L3}}$ (network packet
loss), and $A_{\text{L5}}$ (external SLA delivered by downstream
services running in the same region). Treating these as
independent would under-estimate the risk.

\emph{(2) Weakest-link interpretation.}
When a user observes system availability over a window, they
observe whichever category of failure happened to dominate that
window. A month dominated by deploy mistakes is capped by
$A_{\text{L1}}$; a month dominated by hardware failures is capped
by $A_{\text{L2}}$; a month dominated by downstream SLA violations
is capped by $A_{\text{L5}}$. The observed availability is
approximately the minimum, not the product, of the categorical
ceilings.

\emph{(3) Conservatism and interpretability.}
For the intended use case---estimating a ceiling that practitioners
should not expect to exceed---$\min$ is conservative (it gives the
tighter upper bound among reasonable choices) and interpretable
(the dominating layer is immediately visible). Practitioners can
then ask ``which layer is the weakest link, and how do I raise it?''
rather than reasoning about the compound effect of a product.

\paragraph{Relationship to classical RBD.}
We do not claim that min-composition is universally preferable to
the series-product. For \emph{hardware subsystems within a single
layer}, where component independence is a reasonable approximation
(e.g., disks in a RAID array, replicated VMs in different failure
domains), the classical series-product remains appropriate and is
retained inside the individual layer models below. The contribution
of this section is specifically the \emph{cross-layer} composition:
we compose \emph{layers} with $\min$, while each layer internally
may still use classical composition.

\paragraph{Limitation.}
The min-composition operator is a \emph{modeling choice} motivated
by the weakest-link interpretation, not a theorem derived from
first principles. A practitioner whose system exhibits strong
independence between categories (an unusual situation) should use
the series-product instead. The open-source release allows users
to substitute a different composition operator.

\paragraph{Sensitivity analysis: min vs.\ product.}
To quantify the practical impact of the composition operator
choice, we computed both $\min$ and $\prod$ across the 36~test
topologies. The min-composition yields system availability
0.01--0.10~percentage points higher (more optimistic) than the
series-product (mean: 0.02~pp across the test suite). This gap
grows with the number of layers and with inter-layer availability
differences: when one layer dominates (e.g., $A_{\text{L1}} = 0.998$
while others are $\ge 0.9999$), $\min$ and $\prod$ nearly coincide;
when multiple layers are comparably weak, $\prod$ gives a
substantially lower estimate. Practitioners can compare both operators on their
topology to decide which better matches their operational
experience. The key question is empirical: do observed outages
tend to be dominated by a single category per window
(favoring~$\min$) or by simultaneous contributions from multiple
categories (favoring~$\prod$)?

\paragraph{Direction of the gap.}
The min-composition is more \emph{optimistic} (higher availability estimate)
than the series-product. For a resilience analysis tool, this direction
warrants caution: practitioners who rely on the ceiling estimate may
under-estimate risk relative to the classical approach. We recommend
that practitioners compute both operators on their topology and treat
the series-product as a pessimistic lower bound and the min-composition
as an optimistic upper bound. The gap between the two quantifies the
sensitivity to the independence assumption.

\subsection{Layer 1: Software Limit}
The practical ceiling accounting for deployment downtime, human error,
and configuration drift:
\begin{equation}
  A_{\text{L1}} = 1 - \left(
    \frac{d \cdot \bar{t}_{\text{deploy}}}{T_{\text{month}}}
    + p_{\text{human}} + p_{\text{drift}}
  \right)
\end{equation}
where $d$ is deployments per month, $\bar{t}_{\text{deploy}}$ is mean
deploy downtime (seconds), $T_{\text{month}} = 30.44 \times 86400$~s,
$p_{\text{human}}$ is human error probability per month, and
$p_{\text{drift}}$ is configuration drift probability.
All layer availabilities are clamped to $[0, 1]$:
$A_{\text{Lk}} = \max(0, \min(1, \ldots))$.

\subsection{Layer 2: Hardware Limit}
The physical ceiling from MTBF, MTTR, redundancy, and failover.
When component-specific MTBF/MTTR values are not provided, the
implementation uses configurable defaults (e.g., app server
MTBF~=~2160h, MTTR~=~5min; database MTBF~=~4320h, MTTR~=~30min).
These defaults are \emph{engineering estimates}, not empirically
derived from field data. Practitioners should supply measured
values from their monitoring systems; the defaults are intended
only as starting points for exploratory analysis.
The formulas are:
\begin{equation}
  A_{\text{single}}(c) = \frac{\text{MTBF}(c)}{\text{MTBF}(c) + \text{MTTR}(c)}
\end{equation}
\begin{equation}
  A_{\text{tier}}(c) = 1 - \bigl(1 - A_{\text{single}}(c)\bigr)^{n(c)}
\end{equation}
where $n(c)$ is the number of replicas.  If failover is enabled:
\begin{equation}
  A_{\text{tier}}^{\prime}(c) = A_{\text{tier}}(c) \cdot (1 - f_{\text{fo}}(c))
\end{equation}
where $f_{\text{fo}}(c)$ is the failover downtime fraction, computed from
failover event frequency and promotion time.  The system hardware
availability is the product over critical-path components:
\begin{equation}
  A_{\text{L2}} = \max\!\left(0, \min\!\left(1,\,
    \prod_{c \in C_{\text{crit}}} A_{\text{tier}}^{\prime}(c)\right)\right)
\end{equation}

\subsection{Layer 3: Runtime Noise Floor}
The ceiling imposed by irreducible runtime overhead---network packet
loss, garbage collection pauses, and kernel scheduling jitter---as a
standalone environment-specific bound, independent of the hardware
limit:
\begin{equation}
  A_{\text{L3}} = \max\!\left(0, \min\!\left(1,\,
    (1 - \bar{p}_{\text{loss}}) \cdot (1 - \bar{f}_{\text{gc}})\right)\right)
\end{equation}
where $\bar{p}_{\text{loss}}$ is the mean packet loss rate across
components and $\bar{f}_{\text{gc}}$ is the mean fraction of time
spent in GC pauses:
\begin{equation}
  \bar{f}_{\text{gc}} = \frac{1}{|C|} \sum_{c \in C}
    \frac{\mathit{gc\_pause}(c) \cdot \mathit{gc\_freq}(c)}{1000}
\end{equation}
Note: the earlier formulation $A_{\text{L3}} = A_{\text{L2}} \cdot
(1-\bar{p}_{\text{loss}}) \cdot (1-\bar{f}_{\text{gc}})$ made
$A_{\text{L3}} \le A_{\text{L2}}$ by construction, which caused
$\min(A_{\text{L2}}, A_{\text{L3}}) = A_{\text{L3}}$ always,
rendering~$A_{\text{L2}}$ redundant in the cross-layer min.
The revised formula treats Layer~3 as an independent runtime
noise ceiling.

\subsection{Layer 4: Operational Limit}
Captures the human factor---incident detection and resolution speed:
\begin{equation}
  A_{\text{L4}} = \max\!\left(0, \min\!\left(1,\,
    1 - \frac{n_{\text{inc}} \cdot \bar{t}_{\text{resp}}}{8760 \cdot p_{\text{cov}}}\right)\right)
\end{equation}
where $n_{\text{inc}}$ is incidents per year,
$\bar{t}_{\text{resp}}$ is mean response time (hours),
and $p_{\text{cov}} \in (0, 1]$ is on-call coverage fraction
(1.0 = 24/7, 0.33 = business hours).
The clamp is necessary: for high incident rates or long response times
the formula can yield $A_{\text{L4}} < 0$ without it
(e.g., $n_{\text{inc}} = 500$, $\bar{t}_{\text{resp}} = 20$h,
$p_{\text{cov}} = 0.33$ gives $A_{\text{L4}} \approx -2.46$).
Per-component operational readiness factors (runbook coverage,
automation percentage) further reduce the effective response time.

\subsection{Layer 5: External SLA Cascading}
If the system depends on $m$ external services, each with SLA $a_i$:
\begin{equation}
  A_{\text{L5}} = \max\!\left(0, \min\!\left(1,\,\prod_{i=1}^{m} a_i\right)\right)
\end{equation}
This layer captures the multiplicative compounding of external SLA
dependencies.
When multiple external services are hosted by the same cloud
provider, their SLAs are correlated (a provider-level outage
degrades all). In such cases, the series-product over-estimates
the compounding effect; practitioners should group co-located
services and apply $\min$ within each provider group.

\subsection{Cascade Path Availability}
For a cascade path $P = (c_1, c_2, \ldots, c_k)$ through the
dependency graph, each edge introduces an attenuation factor:
\begin{equation}
  A_{\text{path}}(P) = \prod_{i=1}^{k-1} \alpha(\mathit{dep}(c_i, c_{i+1}))
    \cdot A_{\text{system}}
\end{equation}
where $\alpha(\texttt{requires}) = 1.0$,
$\alpha(\texttt{optional}) \le 0.5$, and
$\alpha(\texttt{async}) \le 0.3$.
These attenuation factors are configurable engineering parameters,
not empirically measured values (see \S\ref{sec:discussion}).
%% ====================================================================
%%  7. ILLUSTRATIVE BACKTEST (NOT PROSPECTIVE EVALUATION)
%% ====================================================================
\section{Illustrative Backtest}\label{sec:eval}

\textbf{Scope and limitations of this section.}
This section does \emph{not} constitute a prospective evaluation.
All results below are \emph{post-hoc}: for each incident, the
topology graph was constructed after reading the public
post-mortem, so the graph already encodes which components were
affected. The $F_1$ scores reported below are therefore
\textbf{trivially achievable by construction}---any correct
graph traversal over a topology that encodes the ground truth
will reproduce the known cascade path. The purpose of this
section is limited to verifying a \emph{necessary sanity check}:
given an accurate topology and the documented root cause, does
the cascade engine's propagation rules reproduce the known
cascade path under the modeled dependency semantics? A cascade
engine that fails this check would be incorrect; passing it is
necessary but not sufficient evidence of utility.

We explicitly do not claim predictive accuracy on unseen
topologies. Whether a formal LTS cascade engine generalizes
to infrastructures whose dependency types and circuit breakers
are unknown at the time of failure remains an open question.
Prospective evaluation on live systems (e.g., against
DeathStarBench traces as in Krasnovsky~\cite{krasnovsky2026})
is needed to establish predictive validity and is left to
future work.

\subsection{Post-Hoc Cascade Path Reproduction}

We assembled 36~incident topologies from real-world cloud
infrastructure incidents spanning 2017--2023, of which 29~achieve
$F_1 = 1.000$ and 7~exhibit imperfect cascade reproduction
(see below). For each incident, we constructed a representative
topology graph from the public post-mortem, injected the documented
root cause as a fault, and compared the predicted cascade path
against the incident report. Because the topology encodes the
answer, the cascade path $F_1$ is expected to be 1.000; any
deviation indicates a modeling limitation rather than a meaningful
predictive result.

\begin{table*}[t]
\centering
\caption{Post-hoc cascade path reproduction on real-world incidents
  (8 of 29 with $F_1 = 1.000$ shown; full results in repository).
  Because topologies were constructed from incident reports, the
  path $F_1 = 1.000$ is \emph{expected by construction} and does
  not indicate predictive power.
  Sev.\ Acc.\ = Severity Accuracy (non-trivial).}
\label{tab:backtest}
\small
\begin{tabular}{@{}llccccc@{}}
\toprule
\textbf{Incident} & \textbf{Year} & \textbf{Pred.\ Cascade} & \textbf{Actual Cascade} & \textbf{Prec.} & \textbf{Rec.} & \textbf{$F_1$} \\
\midrule
AWS us-east-1 outage      & 2021 & 7 components & 7 components & 1.000 & 1.000 & 1.000 \\
AWS S3 key-space           & 2017 & 5 components & 5 components & 1.000 & 1.000 & 1.000 \\
Meta BGP misconfiguration  & 2021 & 4 components & 4 components & 1.000 & 1.000 & 1.000 \\
Cloudflare BGP misconfiguration        & 2022 & 3 components & 3 components & 1.000 & 1.000 & 1.000 \\
GCP networking             & 2019 & 6 components & 6 components & 1.000 & 1.000 & 1.000 \\
Azure AD outage            & 2023 & 5 components & 5 components & 1.000 & 1.000 & 1.000 \\
GitHub Actions outage      & 2022 & 4 components & 4 components & 1.000 & 1.000 & 1.000 \\
Fastly CDN outage          & 2021 & 3 components & 3 components & 1.000 & 1.000 & 1.000 \\
\midrule
\multicolumn{2}{@{}l}{\textbf{Average (29 of 36 incidents)}} & --- & --- & 1.000 & 1.000 & 1.000\textsuperscript{$\dagger$} \\
\multicolumn{2}{@{}l}{\textbf{Severity Accuracy}} & \multicolumn{5}{c}{\textbf{0.819}} \\
\bottomrule
\end{tabular}
\\[2pt]
{\footnotesize \textsuperscript{$\dagger$}Expected by construction;
  see \S\ref{sec:eval} preamble.}
\end{table*}

\subsection{Perturbation Analysis: Is the Result Trivial?}

To demonstrate that the cascade path reproduction is not achieved
by a degenerate engine (e.g., one that simply marks all reachable
nodes as failed), we performed a perturbation analysis. For each
of the 36~incident topologies in our test suite, we systematically
perturbed the topology in three ways and re-ran the cascade engine:

\begin{enumerate}[nosep]
\item \textbf{Edge deletion}: removed one randomly-selected
  dependency edge. Average $F_1$ dropped from 0.971 to 0.818
  (average recall dropped to 0.724 because the engine could no
  longer reach components through the deleted path).
\item \textbf{Edge addition}: added one spurious \texttt{requires}
  dependency from a new component to the root-cause node. Average
  precision dropped to 0.748 (the engine propagated failure to
  the spurious component, which was not in the ground-truth
  cascade path).
\item \textbf{Dependency type swap}: changed one randomly-selected
  \texttt{requires} edge to \texttt{optional}. Average $F_1$
  dropped to 0.959 (moderate effect because optional dependencies
  attenuate to \textsc{Degraded} rather than \textsc{Down},
  partially stopping further cascade).
\end{enumerate}

These perturbations confirm that the engine's behavior is
\emph{sensitive to topology structure}: incorrect topologies
produce measurably worse results. Edge deletion has the largest
impact ($\Delta F_1 = -0.153$), confirming that the engine does
not trivially mark all reachable nodes as failed. This does not
establish predictive power, but it does establish that the engine
is non-degenerate.
A limitation of the edge addition perturbation is that it
introduces a new node not present in the original topology,
which inflates false positives in a somewhat artificial way.
A more realistic perturbation would add a spurious edge between
two existing nodes; we chose the node-addition variant for
simplicity and to avoid confounds from existing dependency
structure.

\paragraph{Baseline imperfection on the full test suite.}
We note that the baseline $F_1$ across all 36~incident topologies
is 0.971, not 1.000. Seven incidents exhibit over-propagation
(6~cases, where the engine predicts cascade to a component not
in the ground truth) or under-propagation (1~case, where the
engine misses an affected component) relative to the
ground-truth cascade path. This reflects a real limitation:
the engine's propagation rules (particularly for optional and
async dependencies) do not perfectly match the actual cascade
behavior in all cases. Table~\ref{tab:backtest} shows 8 representative incidents from
the 29 that achieve $F_1 = 1.000$; the full 36-incident results
are available in the repository.

Table~\ref{tab:imperfect} summarizes the seven imperfect incidents by failure mode and root cause.

\begin{table}[h]
\centering
\caption{Seven incidents with imperfect cascade reproduction ($F_1 < 1.000$).
  All six over-propagation cases arise from optional-dependency attenuation:
  the engine propagates \textsc{Degraded} to components that were
  isolated by circuit breakers or feature flags in the real incident.
  The one under-propagation case reflects a missing implicit dependency
  (shared storage or DNS) not captured in the topology YAML.}
\label{tab:imperfect}
\small
\begin{tabular}{@{}lllp{5.5cm}@{}}
\toprule
\textbf{Incident} & \textbf{Year} & \textbf{Mode} & \textbf{Root Cause of Mismatch} \\
\midrule
Provider A outage  & 2017 & over  & Optional dep.\ propagated to isolated replica group \\
Provider B routing & 2019 & over  & CB trip masked downstream; engine still propagated \\
Provider C auth    & 2020 & over  & Feature-flag disabled path not modeled as isolated \\
Provider D storage & 2021 & over  & Async dep.\ to read-replica treated as active path \\
Provider E BGP     & 2022 & over  & BGP scope limited to 19 DCs; engine applied globally \\
Provider F SSL     & 2022 & over  & Certificate scope mismatch; engine over-generalized \\
Provider G control & 2023 & under & Implicit shared-state dependency absent from YAML \\
\bottomrule
\end{tabular}
\end{table}

A separate evaluation on a 54-incident superset spanning 21 providers
and 2017--2024 shows $F_1 = 0.87$ and recall $= 0.81$, consistent
with the modeling limitations described above; details are reported
in a companion short paper.

\paragraph{Severity accuracy.}
The severity accuracy of 0.819 is the only non-trivial metric in
this section: severity scores are computed by the engine's internal
scoring function (a weighted combination of health-state impact,
cascade spread, and likelihood; 0--10 scale), so they represent
a genuine---if limited---signal. The metric maps the continuous score to four levels
(critical: $\ge 7.0$, major: $\ge 4.0$, minor: $\ge 1.0$,
low: $< 1.0$) and compares against the level reported in
the incident post-mortem. The moderate accuracy
reflects both modeling limitations and the inconsistent severity
definitions used across different cloud providers' incident
reports. The scoring function and thresholds are documented
in the open-source implementation (\texttt{cascade.py}).

\subsection{Scale}
The \textsc{FaultRay} repository contains 512~Python source files
and over 30{,}000~test functions. The core modules relevant to the
contributions of this paper (cascade engine, availability model,
and infrastructure graph model) comprise approximately 12~source
files and 2{,}000~lines of code; the remainder covers domain-specific
extensions, API endpoints, integrations, and tooling that are part
of the open-source release but not claimed as research contributions.
The cascade engine processes a 50-component topology in under 10\,ms
on commodity hardware (Apple M1, 16\,GB RAM).

\subsection{Comparison with Existing Tools}

Table~\ref{tab:comparison} summarizes the structural differences
between \textsc{FaultRay} and representative chaos engineering
platforms. We emphasize that these are different operating
points, not a value judgment: production fault injection and
in-memory analytical simulation address different use cases.

\begin{table}[t]
\centering
\caption{Structural differences between \textsc{FaultRay} and
  representative chaos engineering platforms. These are distinct
  operating points, not absolute rankings.}
\label{tab:comparison}
\small
\begin{tabular}{@{}lcccc@{}}
\toprule
\textbf{Feature} & \textbf{FR} & \textbf{Gr} & \textbf{Sb} & \textbf{FIS} \\
\midrule
Analytical (in-memory) & \checkmark & \texttimes & \texttimes & \texttimes \\
Formal cascade proofs  & \checkmark & \texttimes & \texttimes & \texttimes \\
Cascade propagation    & \checkmark & Partial    & Partial    & \texttimes \\
$N$-layer decomposition & \checkmark & \texttimes & \texttimes & \texttimes \\
Zero production touch  & \checkmark & \texttimes & partial & \texttimes \\
Open source (Apache 2.0) & \checkmark & \texttimes & Partial    & \texttimes \\
Real fault injection  & \texttimes & \checkmark & \checkmark & \checkmark \\
\bottomrule
\end{tabular}

\smallskip
\raggedright
\footnotesize FR = FaultRay, Gr = Gremlin, Sb = Steadybit, FIS = AWS FIS.
The bottom row (real fault injection) is a capability that
\textsc{FaultRay} deliberately does not provide.
\end{table}

\paragraph{Limitations of this comparison.}
Table~\ref{tab:comparison} is a \emph{structural feature
matrix}, not a quantitative accuracy comparison. Direct
accuracy comparison between \textsc{FaultRay} and production
fault injection tools is not possible without a shared benchmark
topology where both approaches are evaluated against the same
ground truth. Krasnovsky~\cite{krasnovsky2026} evaluated against
DeathStarBench (MAE~$\leq$~0.0004); replicating this benchmark
in \textsc{FaultRay} would require manually constructing the
topology from Jaeger traces, which we have not done. We
acknowledge this as an important gap: \textbf{no quantitative
comparison to any baseline exists in this paper.}

\paragraph{Comparison with naive graph reachability.}
For binary cascade path prediction (``is component~$c$ affected
or not?''), the cascade engine produces identical results to
naive reverse-BFS on all 36~test topologies ($F_1 = 0.971$ for
both). This equivalence holds even on topologies containing
\texttt{optional} and \texttt{async} dependencies, because the
binary $F_1$ metric treats \textsc{Degraded} and \textsc{Down}
identically as ``affected.'' The cascade engine's \emph{actual}
added value over BFS lies in severity differentiation
(assigning \textsc{Degraded} vs.\ \textsc{Overloaded} vs.\
\textsc{Down} based on dependency type and weight), latency
cascade simulation, and the formal guarantees
(Theorems~\ref{thm:termination}--\ref{thm:blast}), none of
which are captured by the binary $F_1$ metric.

\paragraph{Min-composition vs.\ series-product.}
On all 36~test topologies, the 3-layer min-composition yields
system availability estimates 0.01--0.10~percentage points
\emph{higher} (more optimistic) than the classical
series-product (mean: 0.02~pp), corresponding to
50--540~fewer minutes of estimated annual downtime
(mean: 103~minutes). This is consistent with the
intended interpretation: $\min$ produces a \emph{ceiling}
(what you cannot exceed), while the series-product produces
a \emph{point estimate} (what you expect under independence).
Neither is universally correct; the choice depends on whether
the practitioner's layers exhibit correlated or independent
failure modes. The open-source release allows users to
substitute a different composition operator.

%% ====================================================================
%%  8. IMPLEMENTATION
%% ====================================================================
\section{Implementation}\label{sec:impl}

\textsc{FaultRay} is implemented in Python 3.11+ and uses NetworkX~\cite{networkx}
for graph operations, Pydantic for schema validation, and FastAPI for
the REST API.  The software is distributed as an open-source package on
PyPI under the Apache License 2.0.

\textbf{Infrastructure graph model.}
Topologies are defined in JSON or YAML.  Each component specifies type,
capacity, resource metrics, and redundancy configuration.
Dependencies specify type
(\texttt{requires}/\texttt{optional}/\texttt{async}), weight, latency,
circuit breaker configuration, and retry strategy.

\textbf{CLI interface.}
The \texttt{faultray} CLI provides commands for topology validation,
single-fault simulation, traffic spike analysis, latency cascade
modeling, availability limit computation, and backtest execution against
historical incidents.

\textbf{API and integrations.}
A FastAPI server exposes the simulation engines via REST endpoints.
Integrations with Terraform (infrastructure-as-code scanning),
Prometheus (live metrics ingestion), and CI/CD pipelines (quality gates)
are provided.

\textbf{Extensibility.}
New fault types, component types, and simulation engines can be
added via a plugin architecture.  Custom scoring functions allow
organizations to weight severity according to their specific SLO
definitions.

\textbf{Reproducibility.}
All 36~incident topology files used in the backtest are included
in the repository at \texttt{tests/incidents/*.yaml}. The
perturbation analysis uses a fixed random seed (42) for
reproducibility. The Python version is 3.11+; exact dependency
versions are pinned in \texttt{pyproject.toml}.

%% ====================================================================
%%  9. DISCUSSION
%% ====================================================================
\section{Discussion and Limitations}\label{sec:discussion}

We describe the limitations of this work in detail because the
intended contribution of the paper is enablement of a research
prototype, not an empirical novelty claim. Practitioners considering
\textsc{FaultRay} should understand what it does and does not
establish.

\textbf{No prospective validation.}
All backtest results in \S\ref{sec:eval} are post-hoc: topologies
were constructed from incident reports, so cascade path $F_1 = 1.000$
is expected by construction rather than earned by prediction. The
perturbation analysis (\S\ref{sec:eval}) demonstrates non-degeneracy
but not predictive power. We explicitly do not claim predictive
accuracy on unseen infrastructures. Prospective evaluation---ideally
on trace-derived topologies such as DeathStarBench, where the engine
must predict cascade paths from topology alone without foreknowledge
of the outcome---remains the most important open validation task.

\textbf{Relationship to concurrent work (Krasnovsky 2026).}
Krasnovsky's ICSE-NIER~2026 paper~\cite{krasnovsky2026} proposes
in-memory graph simulation with Monte-Carlo sampling, evaluated on
DeathStarBench with quantitative mean absolute error bounds. That
paper provides empirical results that this paper does not, and we
recommend it for readers interested in trace-driven graph discovery.
The two papers differ in method (Monte-Carlo vs.\ formal LTS) and
scope (single-layer vs.\ multi-layer availability). We view them as
complementary operating points in the design space of in-memory
chaos analysis.

\textbf{Proof sketches, not mechanized proofs.}
The correctness properties in \S\ref{sec:cascade} (Theorems~1--7)
are accompanied by informal proof sketches, not proofs verified
in a proof assistant such as Coq, Isabelle/HOL, or TLA+. While
we have attempted to make the reasoning rigorous (e.g., explicit
progress measures for termination, case analysis for monotonicity),
informal arguments can contain subtle errors that mechanical
verification would catch. We consider mechanized verification of
the core properties an important future work item.

\textbf{Modeling simplifications.}
\textsc{FaultRay} operates on a static topology model and does not
capture runtime dynamics such as auto-healing, load balancer failover
timing, auto-scaling reaction time, or Kubernetes pod rescheduling
latency. Dependency types (\texttt{requires} / \texttt{optional} /
\texttt{async}) are a simplification of the far richer set of
coupling semantics present in real distributed systems.

\textbf{Coarse health status granularity.}
The four-level health model (\textsc{Healthy}, \textsc{Degraded},
\textsc{Overloaded}, \textsc{Down}) does not capture gradient
degradation within a level. For example, \textsc{Degraded} at
2$\times$ latency and \textsc{Degraded} at 10$\times$ latency are
treated identically by the cascade propagation rules, yet their
impact on user experience and downstream systems may differ
substantially. Similarly, \textsc{Overloaded} at 71\% CPU
utilization and 99\% are conflated. A continuous health metric
would enable finer-grained propagation modeling but would
complicate the formal properties (e.g., monotonicity over a
partial order becomes monotonicity over a total order on
$[0, 1]$). We chose discrete levels for simplicity and
interpretability; extending to continuous health is left to
future work.

\textbf{Min-composition is a modeling choice, not a theorem.}
As discussed in \S\ref{sec:availability}, the decision to compose
categorical availability layers with $\min$ rather than the classical
series-product is motivated by a weakest-link argument and by the
observation that cloud infrastructure layers fail in correlated
fashion. It is not derived from first principles. Practitioners whose
system genuinely exhibits independence between categorical layers
should use a different composition operator.

\textbf{Scope of this paper.}
An earlier version of this manuscript (v11) presented an AI-agent
cross-layer failure model, a 10-mode taxonomy of agent failures,
and a hallucination probability function $H(a, D, I)$ with
hardcoded parameters. We have removed that section because
subsequent agent failure taxonomies---notably
MAST~\cite{cemri2025mast}, Microsoft's Taxonomy of Failure Mode in
Agentic AI Systems, and work on AgentOps---cover the taxonomy
contribution with substantially more empirical rigor (1600+
annotated traces, inter-annotator agreement $\kappa = 0.88$).
A broader framework for AI agent reliability engineering has also
been proposed~\cite{airel2026}.
The
hallucination parameters $\delta = 0.5$, $\omega = 0.3$ in the v11
paper were engineering judgment and not measured. We therefore
scope this v12 paper to the infrastructure cascade engine and the
$N$-layer availability model, and defer empirical AI-agent failure
modeling to domain specialists. The open-source implementation
retains the AI agent features as a tool capability, but they are
not claimed as a research contribution in this paper.

\textbf{Not a compliance tool.}
\textsc{FaultRay} is a research prototype. It is not a certified
compliance tool under DORA, FISC, SOC~2, ISO~27001, or any other
regulatory regime. Its outputs must not be used as audit evidence
without independent legal and technical review. Regulated entities
should use certified tooling operated by licensed assessors for
compliance purposes.

\textbf{Patent status.}
A US provisional patent application covering the cascade engine and
$N$-layer availability decomposition has been filed with the USPTO
(Application No.~64/010,200, filed 2026-03-19). The release of this
paper under an open preprint, together with the Apache~License~2.0
open-source release, is intended as defensive publication in
addition to the provisional filing. Note that the Apache~License~2.0
(\S3) includes an express patent grant to downstream users of the
licensed software; the provisional filing preserves rights for
uses outside the scope of the Apache license (e.g., independent
reimplementation).

\textbf{Scalability.}
The $O(|V| + |E|)$ cascade engine scales to hundreds of components
in tens of milliseconds on commodity hardware. Very large
topologies ($>$10{,}000 components) may require partitioning
strategies. The per-scenario simulation is embarrassingly parallel.

\textbf{What this paper does not claim.}
To be explicit, this paper does \emph{not} claim:
(i)~predictive accuracy on unseen topologies;
(ii)~superiority over production fault injection for use cases
requiring empirical fault discovery;
(iii)~novelty over every prior work on cascade failure modeling; or
(iv)~that $\min$-composition is universally preferable to the
series-product.
\textbf{Threats to validity.}
\emph{Internal:} The proof sketches are not mechanically verified;
subtle errors may exist. The severity accuracy metric (0.819)
depends on manual mapping of incident reports to severity levels,
which involves subjective judgment.
\emph{External:} The 36~test topologies are simplified
representations of real incidents (3--7 components each), not
full production topologies (which may have hundreds of components
with complex dependency types). Results may not generalize to
larger or more complex topologies.
\emph{Construct:} The binary $F_1$ metric does not capture
severity differentiation, which is the cascade engine's primary
added value over naive graph reachability.

What the paper \emph{does} claim is that the specific combination of
(I)~formal LTS cascade semantics with termination and complexity
arguments (proof sketches), (II)~cross-layer $\min$-composition of categorical
availability ceilings, and (III)~an Apache~2.0 in-memory analytical
implementation, constitutes a useful operating point that, to the
best of our knowledge, had not been disclosed in this combined form
as of the provisional filing date (2026-03-19).

%% ====================================================================
%%  10. CONCLUSION
%% ====================================================================
\section{Conclusion}\label{sec:conclusion}

We presented \textsc{FaultRay}, a research prototype for in-memory
infrastructure resilience analysis. The paper contributes two core
technical elements:

\begin{enumerate}[nosep]
\item A \textbf{cascade propagation engine} formalized as a Labeled
  Transition System over a dependency graph, with termination
  on finite graphs, $O(|V| + |E|)$ time complexity, monotonicity,
  causality, and bounded blast radius (proof sketches in
  \S\ref{sec:cascade}; not mechanically verified).
\item An \textbf{$N$-layer availability limit model} that composes
  categorical availability ceilings (software, hardware, runtime
  noise floor, operational, external SLA) with the $\min$ operator rather than
  the classical series-product, motivated by a weakest-link argument
  appropriate for cloud infrastructure where categorical layers
  exhibit correlated failure modes (\S\ref{sec:availability}).
\end{enumerate}

A post-hoc backtest on 36 historical cloud incidents (29 with
$F_1 = 1.000$, 7 with modeling imperfections) confirms that the
cascade engine largely reproduces known cascade paths when
topologies encode the ground truth (a necessary sanity check,
not a predictive result); perturbation analysis shows
non-degenerate sensitivity to topology structure. Prospective
evaluation remains open. Concurrent work by Krasnovsky~\cite{krasnovsky2026}
(ICSE-NIER~2026) addresses the same operating point with Monte-Carlo
sampling, and we view the two approaches as complementary.

\textsc{FaultRay} is released as open source under Apache~License~2.0
at \url{https://github.com/mattyopon/faultray} and
\url{https://pypi.org/project/faultray/}.
A US provisional patent application covering the cascade engine and
$N$-layer availability decomposition has been filed with the USPTO
(Application No.~64/010,200, 2026-03-19).

This is a research prototype, not a certified compliance tool. Its
outputs must not be used as audit evidence without independent legal
and technical review.

%% ====================================================================
%%  REFERENCES
%% ====================================================================
\section*{Acknowledgments}
The author thanks the open-source community for feedback on early versions of FaultRay.
Claude (Anthropic) was used for drafting assistance in preparing this manuscript;
the author assumes full responsibility for the accuracy of all technical content.

\begin{thebibliography}{99}

\bibitem{basiri2016}
A.~Basiri, N.~Behnam, R.~de~Rooij, L.~Hochstein, L.~Kosewski,
J.~Reynolds, and C.~Rosenthal,
``Chaos engineering,''
\textit{IEEE Software}, vol.~33, no.~3, pp.~35--41, 2016.

\bibitem{gremlin2020}
Gremlin, Inc.,
``Gremlin: Chaos engineering platform,''
2020. [Online]. Available: \url{https://www.gremlin.com/}

\bibitem{steadybit2023}
Steadybit GmbH,
``Steadybit: Chaos engineering for Kubernetes,''
2023. [Online]. Available: \url{https://www.steadybit.com/}

\bibitem{aws-fis}
Amazon Web Services,
``AWS Fault Injection Simulator,''
2021. [Online]. Available:
\url{https://aws.amazon.com/fis/}

\bibitem{litmuschaos}
LitmusChaos,
``LitmusChaos: Cloud-native chaos engineering,''
CNCF Incubating Project, 2020. [Online]. Available:
\url{https://litmuschaos.io/}

\bibitem{hpe-spof}
Hewlett Packard Enterprise,
``Detecting single points of failure,''
U.S. Patent 9,280,409 B2, Mar.~8, 2016.

\bibitem{keller1976}
R.~M.~Keller,
``Formal verification of parallel programs,''
\textit{Communications of the ACM}, vol.~19, no.~7,
pp.~371--384, 1976.

\bibitem{trivedi2001}
K.~S.~Trivedi,
\textit{Probability and Statistics with Reliability, Queuing, and
Computer Science Applications}, 2nd~ed.
New York: Wiley, 2001.

\bibitem{networkx}
A.~A.~Hagberg, D.~A.~Schult, and P.~J.~Swart,
``Exploring network structure, dynamics, and function using NetworkX,''
in \textit{Proc.\ SciPy}, pp.~11--15, 2008.

\bibitem{meta2021}
Janardhan, S.,
``More details about the October 4 outage,''
Meta Engineering Blog, Oct.~2021. [Online]. Available:
\url{https://engineering.fb.com/2021/10/05/networking-traffic/outage-details/}

\bibitem{aws-s3-2017}
Amazon Web Services,
``Summary of the Amazon S3 service disruption in the Northern
Virginia (US-EAST-1) region,''
Feb.~2017. [Online]. Available:
\url{https://aws.amazon.com/message/41926/}

\bibitem{cloudflare2022}
Cloudflare, Inc.,
``Cloudflare outage on June 21, 2022,''
Cloudflare Blog, Jun.~2022. [Online]. Available:
\url{https://blog.cloudflare.com/cloudflare-outage-on-june-21-2022/}

\bibitem{krasnovsky2026}
A.~A.~Krasnovsky,
``Model Discovery and Graph Simulation: A Lightweight Gateway to
Chaos Engineering,''
in \textit{Proceedings of the 48th International Conference on
Software Engineering, NIER Track (ICSE-NIER~2026)}, Apr.~2026.
arXiv:2506.11176.
\newblock \textit{Distinguished Paper Award.}

\bibitem{casformer2025}
C.~S.-H.~Hsu et al.,
``Cascaded Transformer for robust and scalable SLA decomposition
via amortized optimization,'' arXiv:2601.11859, 2026.

\bibitem{nasa-fta}
W.~Vesely, M.~Stamatelatos, J.~Dugan, J.~Fragola, J.~Minarick, and
J.~Railsback,
\textit{Fault Tree Handbook with Aerospace Applications},
NASA Office of Safety and Mission Assurance, 2002.

\bibitem{iec61025}
International Electrotechnical Commission,
\textit{IEC 61025: Fault Tree Analysis (FTA)},
2nd ed., 2006.

\bibitem{arp4761}
SAE International,
\textit{ARP4761: Guidelines and Methods for Conducting the Safety
Assessment Process on Civil Airborne Systems and Equipment},
1996.

\bibitem{iii-model-2025}
Y.~Tang, J.~Piao, H.~Wang, S.~Rajib, and Y.~Li,
``Predicting cascade failures in interdependent urban
infrastructure networks,'' arXiv:2503.02890, 2025.

\bibitem{pnas-cascade-2019}
J.~Duan et al.,
``Universal behavior of cascading failures in interdependent
networks,''
\textit{Proceedings of the National Academy of Sciences}, 2019.

\bibitem{sla-survey-2024}
``Service Level Agreement in cloud computing: Taxonomy, prospects,
and challenges,'' \textit{Internet of Things}, vol.~25, 101126,
2024. doi:10.1016/j.iot.2024.101126.

\bibitem{cemri2025mast}
M.~Cemri et al.,
``MAST: A Multi-Agent System Failure Taxonomy,''
arXiv:2503.13657, to appear in NeurIPS 2026.

\bibitem{chaostwin2021}
P.~Pecka, T.~Rosa, K.~Zimmermann, and V.~Petrik,
``ChaosTwin: A chaos engineering and digital twin approach for
the validation of virtualized network functions,''
in \textit{Proc.\ CNSM}, 2021.

\bibitem{farokhi2024}
S.~Farokhi et al.,
``Chaos engineering for resilience assessment of digital twins,''
\textit{IEEE Transactions on Industrial Informatics}, vol.~20,
no.~1, pp.~1134--1143, 2024.

\bibitem{chaos-mlr-2025}
O.~S.~Joshua et al.,
``Chaos engineering: A multi-vocal literature review,''
\textit{ACM Computing Surveys}, vol.~58, no.~7, 2025.
doi:10.1145/3777375.

\bibitem{vlsts2022}
V.~Zamfir, M.~Calancea, D.~Diaconescu, W.~Chmielarz,
and J.~Tu\v{s}il,
``Validating labelled state transition and message production
systems,'' arXiv:2202.12662, 2022.

\bibitem{buldyrev2010}
S.~V.~Buldyrev, R.~Parshani, G.~Paul, H.~E.~Stanley, and
S.~Havlin,
``Catastrophic cascade of failures in interdependent networks,''
\textit{Nature}, vol.~464, pp.~1025--1028, 2010.

\bibitem{avizienis2004}
A.~Avizienis, J.-C.~Laprie, B.~Randell, and C.~Landwehr,
``Basic concepts and taxonomy of dependable and secure computing,''
\textit{IEEE Transactions on Dependable and Secure Computing},
vol.~1, no.~1, pp.~11--33, 2004.

\bibitem{nasar2026}
M.~Nasar,
``Formal validation for microgrid reliability based on TLA+,''
\textit{Journal of Electrical Engineering}, vol.~77, no.~1,
pp.~46--56, 2026. doi:10.2478/jee-2026-0006.

\bibitem{airel2026}
S.~Rabanser, S.~Kapoor, P.~Kirgis, K.~Liu, S.~Utpala,
and A.~Narayanan,
``Towards a Science of AI Agent Reliability,''
arXiv:2602.16666, Feb.~2026.

\end{thebibliography}

\end{document}
